% !TEX root = /home/eb/.vscode/extensions/qftrules.programme-colle/tmp/Exercice.tex
%%%%%%%%%%%%%%%%%%%%%%%%%%%%%%%%%%%%%%%%%%%%%%%%%%%%%%%

%%%%%%%%%%%%%%%%%%%%%%%%%%%%%%%%%%%%%%%%%%%%%%%%%%%%%%%
\begin{exo}[PC][1][TD]{Manipulation des équations de Maxwell}
  \Notions{\item Équation de Maxwell}
On suppose que règne dans l'espace le champ électromagnétique 
\begin{equation}
  \E(M, t)=f(z) e^{-t / \tau} \ex, \quad \text{ et } \quad \B(M, t)=g(z) e^{-t / \tau} \ey,
\end{equation}
% $$
% \left\{\begin{array}{l}
% \vec{E}(M, t)=f(z) \mathrm{e}^{-t / \tau} \ex \\
% \vec{B}(M, t)=g(z) \mathrm{e}^{-t / \tau} \ey
% \end{array}\right.
% $$
où $\tau$ est une constante de temps et $f, g$ deux fonctions. L'espace est vide de charges et de courants.
\begin{quest}
  \item   Vérifier que ces champs sont compatibles avec les équations de Maxwell-Gauss et Maxwell-Thomson.
  \sol{
  Maxwell-Gauss : $\dive \E = \dr{E_x}{x} = 0$ car $E_x$ indépendante de $x$. De même, Maxwell-Thomson : $\dive \B = \dr{B_y}{y} = 0$ car $B_y$ indépendante de $y$.
  }
  \item  Montrer que l'équation de Maxwell-Faraday impose une relation entre $g$ et $f^{\prime}$.  Montrer que l'équation de Maxwell-Ampère impose une relation entre $f$ et $g^{\prime}$.
  \sol{
  Maxwell-Faraday : $\rot \E = - \dr{\B}{t} \Rar f'(z) \mathrm{e}^{-t / \tau} \ey = \frac{g(z)}{\tau} \mathrm{e}^{-t / \tau} \ey$ ce qui est possible si
  $$
f'(z) = \frac{g(z)}{\tau}.
  $$
  Maxwell-Ampère : $\rot \B = \frac{1}{c^2}\dr{\E}{t} \Rar -g'(z) \mathrm{e}^{-t / \tau} \ex = -\frac{f(z)}{\tau c^2} \mathrm{e}^{-t / \tau} \ex$ ce qui est possible si
  $$
g'(z) = \frac{f(z)}{\tau c^2}.
  $$
  }
  \item En déduire une équation différentielle vérifiée par $f$. La résoudre en prenant comme condition aux limites $f(z=0)=E_{0}$, et en supposant que le champ électrique est nul en $z \rightarrow+\infty$ à tout instant.  En déduire l'expression complète du champ électromagnétique $(\E,\B)$.
  \sol{
  On trouve $f''(z) = \frac{g'(z)}{\tau} = \frac{f(z)}{\tau^2 c^2}$. On pose $\lambda=\tau c$ une distance caractéristique. Alors $f(z) = A e^{-z/\lambda} +  B e^{z/\lambda} = E_0 e^{-z/\lambda} $ d'après les CL.
  }
\end{quest}
\end{exo}

%%%%%%%%%%%%%%%%%%%%%%%%%%%%%%%%%%%%%%%%%%%%%%%%%%%%%%%
\begin{exo}[PC][2][devoir]{Émission radioactive}
  Une masse radioactive, ponctuelle, initialement neutre, située au point $O$ émet, à partir de l’instant
  $t = 0$, des particules $\alpha$ avec une vitesse supposée constante et de façon isotrope : à l’instant $t$ la charge électrique située en $O$ est $q(t) = q_0 (e^{-t/\tau}-1).$
\begin{quest}
\item Calculer le champ électrique $\E(M,t)$ et le champ magnétique $\B(M,t)$ pour $t > 0$ et commenter. On supposera que les symétries de $\E$ sont les mêmes qu’en statique.
\item Exprimer la densité volumique de charge $\rho(M,t)$ et la densité volumique de courant $\j(M,t)$ pour
$t > 0$.
\item Vérifier la compatibilité des résultats obtenus avec la relation locale de conservation de la charge.
\item En déduire la densité volumique d’énergie électromagnétique, le vecteur de Poynting et la puissance volumique fournie par le champ électromagnétique aux particules $\alpha$. Commenter.
\end{quest}
\end{exo}
%%%%%%%%%%%%%%%%%%%%%%%%%%%%%%%%%%%%%%%%%%%%%%%%%%%%%%%


%%%%%%%%%%%%%%%%%%%%%%%%%%%%%%%%%%%%%%%%%%%%
\begin{exo}[PC][2][TD]{Effet Meissner}
  \Notions{\item Équations de Maxwell \task Densité de courant \task Force et énergie magnétique}

Un matériau supraconducteur conduit le courant électrique sans aucune résistance. Les supraconducteurs présentent la propriété de repousser les lignes de champ magnétiques, ce qui permet à un aimant de léviter au-dessus d'un supraconducteur. Cet effet est appelé effet Meissner.

\begin{fig}[]
  \centering
  \begin{subfig}[test][0.4]
  \includegraphics[width=\textwidth]{supra_meissner.jpg}
\end{subfig}
\begin{subfig}[test][0.35]
  \includegraphics[width=\textwidth]{supra.pdf}
\end{subfig}
\capt{(a) Lévitation d'un aimant au-dessus d'un supraconducteur par effet Meissner. (b) Shéma du problème.}
\end{fig}

Dans la \emph{théorie de London} de la supraconductivité, proposé en 1935, le vecteur densité volumique de courant dans le supraconducteur $\j$ vérifie l'équation de London,
\begin{equation}
  \label{eq:london}
  \rot{\j}= - \frac{ne^2}{m} \B,
\end{equation}
où $e$ désigne la charge élémentaire, $n$ la densité des porteurs de charge supraconducteurs, $m$ la masse de l'électron, et $\vv{B}$ désigne le champ magnétique dans le supraconducteur. On impose un champ magnétostatique uniforme $B_0\, \ez$ dans le vide $x<0$. On cherche le champ magnétique dans le supraconducteur, qui occupe l'espace $x>0$, sous la forme stationnaire $\B(M) = B(x) \ez$.

\begin{quest}
% \item Par analyse dimensionnelle, déterminer la dimension physique de $\alpha$.
\item Établir l'équation vérifiée par le champ $\vv{B}(M)$ dans le supraconducteur, et introduire une longueur de London $\lambda$ en fonction de $n$, $e$, et $m$. Résoudre l'équation en utilisant les conditions aux limites.
% , et interpréter
\item Proposer une interprétation physique de la grandeur $\lambda$. Pour l'aluminium, $\lambda=\SI{1,6e-8}{SI}$. Calculer $B(x)$ à une distance $x=\SI{0,1}{\micro\meter}$ sous la surface du supraconducteur pour $B_0=\SI{1}{T}$. Conclure.
\item Calculer la densité de courant $\j(M)=j(x) \ey$ à l'intérieur du supraconducteur. Estimer l'intensité du courant $I$ parcourant un supraconducteur de hauteur $h = \SI{1}{cm}$ soumis à un champ $B_0 = \SI{1}{T}$.
\item Montrer que la distribution de courant exerce une force magnétique sur la surface du supraconducteur, que l'on caractérisera par sa direction, son sens, et la pression correspondante. 
% \textit{Indication:} La force volumique de Laplace exercée par le champ sur le conducteur s'écrit $\frac{\der\vv{f_{\rm{vol}}}}{\der V}=\vv{j}\pv\vv{B}$.
\end{quest}
\end{exo}
%%%%%%%%%%%%%%%%%%%%%%%%%%%%%%%%%%%%%%%%%%%%



%%%%%%%%%%%%%%%%%%%%%%%%%%%%%%%%%%%%%%%%%%%%%%%%%%%%%%%
\begin{exo}[PC][3][TD]{Solénoïde idéal dans l'ARQS magnétique}
  \Notions{\item ARQS \task Solénoïde \task Équations de Maxwell \task Vecteur de Poynting}
Une bobine, de longueur $\ell=\SI{10}{cm}$, de rayon $a=\SI{10}{cm}$ et d'axe $(Oz)$ est constituée d'un enroulement de $n=\SI{5000}{m^{-1}}$ spires circulaires par unité de longueur. Dans toute la suite on se place dans l'ARQS magnétique et on négligera les effets de bord pour l'expression des champs.
\begin{fig}[label]
  \centering
  \includegraphics[width=0.4\columnwidth]{solenoide}
  % \capt{}
\end{fig}
\begin{quest}
  \item Que vaut le champ magnétique $\B$ engendré par la bobine parcourue par un courant d'intensité $I$ ?
  \sol{
Cours : à l’intérieur, \eqbox{\B(M) = \mu_0 n  I \ez.}
À l’extérieur, $\B=\Vec{0}$. Il faut savoir faire la démo.
}
\item Quel vaut l'énergie magnétique $\Etot_{m}$ associée à la bobine ? En déduire la valeur de l'inductance $L$.
\sol{
Deux méthodes :
\begin{liste}
  \item  Énergie : on a $\mc{E}_{\rm{mag}} = \frac{B^2}{2\mu_0} \times (\ell \pi a^2) = \frac{1}{2} \mu_0 \pi n^2a^2 \ell I^2  =   \frac{1}{2}L I^2$ avec
  \eqbox{L =  \mu_0 \pi n^2a^2 \ell.}
  \item Flux : on a $\Phi  = B \times (N \pi a^2) = L I$.
\end{liste}

}
\item La bobine est placée à l'instant $t=0$  en série avec un générateur de tension $e$, de résistance interne $R$ grande devant celle de l'enroulement. Quelle est la loi d'évolution de l'intensité $I(t)$ dans le circuit ?
\sol{
Il s’agit d’un circuit $RL$ série classique. Tension aux bornes de la bobine : $u=L\dd{I}{t} = e-RI$ donc d’après la loi des mailles
$$\dd{I}{t} + \frac{R}{L} I = \frac{e}{L}.$$.
On définit le temps caractéristique $\tau=\frac{L}{R}$ donc $I(t) = I_0 + A e^{-t/\tau})$ avec $I_0 = \frac{e}{R}$. Condition initiale : continuité de $I$ traversant une bobine en $t=0$, et donc $I(0) = I_0+A = 0$
soit $A=-I_0$. Ainsi,
\eqbox{I(t) = I_0 (1-e^{-t/\tau}).}
}
\item En déduire que le champ magnétique engendré par la bobine à l'instant $t$ en tout point à l'intérieur de la bobine est de la forme 
\begin{equation}
  \B(t)=\alpha(1-e^{-\beta t})\ez,
\end{equation}
 où l'on précisera les expressions de $\alpha$ et $\beta$. Que dire à l'extérieur de la bobine ? 

\item Déterminer l'expression du champ électrique $\vv{E}(M,t)$ en tout point $\rm M$, à l'intérieur et à l'extérieur de la bobine, à partir d'une équation de Maxwell.
\sol{
On utilise la formule précédente. On a
$$\B = \mu_0 n I \ez = \mu_0 n I_0 (1-e^{-t/\tau})\ez,$$ donc $\alpha = \mu_0 n I_0 = \frac{\mu_0 n e}{R}$ et $\beta = \frac{1}{\tau} = \frac{R}{L} = \frac{R}{ \mu_0 \pi n^2a^2 \ell}$. À l'extérieur, le champ magnétique est nul car on néglige les effets de bord et on se place dans l'ARQS.

Pour le champ électrique, on utilise la loi de \imp{Maxwell-Faraday},
$$
\rot E = -\dr{\B}{t} = -\frac{\mu_0 n I_0}{\tau}e^{-t/\tau} \ez.
$$
Pour avoir un rotationnel suivant $\ez$, la solution la plus simple est que le champ $\E$ dépende de $r$ et s’enroule suivant $\et$. On trouve alors à partir d’un formulaire :
$$
\frac{1}{r}\dr{rE_{\theta}}{r} = -\frac{\mu_0 n I_0}{\tau}e^{-t/\tau} ,
$$
Soit
\eqbox{\E = -\frac{\mu_0 n I_0 r}{2\tau}e^{-t/\tau} \et.}

\underline{Remarque :} l’équation de Maxwell-Ampère $\rot \B = \mu_0 \epz \dr{\E}{t}$ ne permet pas de déterminer $\E$ puisque à l’ordre dominant $\B$ est uniforme et donc on trouverait $\E = \Vec{0}$ ! Les calculs effectués ici sont approchées, on calcule $\B$ à l’ordre 0 puis $\E$ à l’ordre 1 en $ac/\tau$, dans ce sens.
}
\item Comparer alors les ordres de grandeur des densités volumiques d'énergie magnétique et électrique en $r=a$ en prenant des valeurs numériques raisonnables pour les variables inconnues.
\sol{
Magnétique : $e_M =  \frac{B^2}{2\mu_0} \sim \left( \frac{\mu_0 n e}{R} \right)^2 /2\mu_0 \sim \frac{1}{2} \mu_0 n ^2 e^2/R^2$. \\

Électrique : $e_E =  \epz \frac{E^2}{2} \sim \epz \left( \frac{\mu_0 n I_0 a}{2\tau}\right)^2 /2 \sim \frac{\epz \mu_0^2 n^2 a^2}{2} (e/L)^2 \sim \frac{\epz e^2}{2\pi^2 n^2 a^2 \ell^2}$.
}
\item Déterminer le vecteur de Poynting $\Poy$ ainsi que son flux à travers la surface de la bobine. Interpréter.
\sol{
$\Poy = \frac{\E \pv \B}{\mu_0} = \frac{E_\theta B_z}{\mu_0} \er = -\frac{\mu_0^2 n^2 I_0^2 r}{2\tau}e^{-t/\tau}(1-e^{-t/\tau}) \er$.
Le vecteur de Poynting est radial donc l’énergie est transféré du bobinage vers le centre de la bobine (pour une décharge). On remarque au passage que
$$
\Poy  = -n I \times \mu_0 n \frac{r}{2}\dd{I}{t} \er = -\frac{1}{2} \mu_0 n^2 r I \dd{I}{t}\er.
$$
La puissance totale transférée vers le centre de la bobine vaut (avec $r=a$)
$$
\mc{P} = \iint_S \Poy  \ps \dS = \frac{1}{2} \mu_0 n a I \dd{I}{t} \times 2 \pi a \ell = (\pi \mu_0 n^2 a^2 \ell)  I \dd{I}{t} =   L I \dd{I}{t}.
$$
On a donc conservation de l’énergie
\eqbox{\mc{P} = \dd{}{t}\left( \frac{1}{2} L I^2\right).}
}
\end{quest}

\textsc{Données :} Opérateur rotationnel en coordonnées cylindriques

 \begin{equation*}
  \rot{A} =  \left[\frac{1}{r} \frac{\partial A_{z}}{\partial \theta} - \frac{\partial A_{\theta}}{\partial z}\right]   \vec{e_r} +
				   \left[\frac{\partial A_{r}}{\partial z} - \frac{\partial A_{z}}{\partial r}\right]   \vec{e_\theta} +
				   \left[\frac{1}{r} \frac{\partial (rA_{\theta})}{\partial r} - \frac{1}{r} \frac{\partial A_{r}}{\partial \theta}\right]   \vec{e_z}
 \end{equation*}
\end{exo}
%%%%%%%%%%%%%%%%%%%%%%%%%%%%%%%%%%%%%%%%%%%%%%%%%%%%%%%


\begin{exo}[PC][3][TD]{Condensateur plan idéal en régime variable}
  \Notions{\item ARQS \task Équations de Maxwell \task Théorèmes intégraux \task Géométrie cylindrique}
  On considère un condensateur plan dont les armatures sont deux disques de rayon $a$, de section $S=\pi a^2$, séparées d'une distance $d$, coaxiaux d’axe $(Oz)$. On néglige les effets de bord en supposant $d\ll a$. Les armatures portent des charges électriques $\pm Q(t)$ dépendantes du temps et réparties uniformément sur la surface des armatures.
   % $u(t)=u_0\cos(\omega t)$.
% \begin{figure}[h!]
% \centering
% \begin{center}
% \includegraphics[width=0.55\textwidth]{condo.pdf}
% \end{center}
% \end{figure}
\begin{fig}[label]
  \centering
  \includegraphics[width=0.2\columnwidth]{disk_capa.pdf}
  %\capt{}
\end{fig}

\begin{quest}
\item Rappeler le cadre de l'ARQS magnétique. Que vaudrait le champ magnétique $\Vec{B}_0$ dans cette approximation ? À partir d'une analyse des invariances et symétries du problème, et du théorème de Gauss, montrer que le champ électrique a pour expression 
$$
\Vec{E}_0(M,t) = \frac{Q(t)}{S\epz} \ez.
$$
\sol{
  \begin{liste}
    \item
  Dans l'AQRS magnétique, l'équation de Maxwell-Ampère se simplifie en $\rot \B_0 = \mu_0 \j = \Vec{0}$ dans tout l'espace. On a aussi $\dive \B_0 =0$. Comme le champ magnétique est nul à l'infini, il est nul partout par unicité de la solution : $\Vec{B}_0 = \Vec{0}$.
\item
Les armatures sont planes et infinies : la distribution de charge est invariante pas translation dans le plane ($x,y)$ et donc $\E_0(M,t) = \E(z,t)$. Les plans $(M,\ex,\ez)$ et $(M,\ey,\ez)$ sont des plans de symétrie, et donc $\E_0(M,t) = E_z(z,t)\ez$. On prend pour surface de Gauss un cylindre d'axe $(Oz)$ de section $s$ quelconque dont une extrémité est à la hauteur $z$ et comprise entre les armatures, et l'autre extrémité est en dehors des armatures. Comme le champ est nul en dehors des armatures, le théorème de Gauss donne alors 
$$
\oiint \E_0(M,t) \ps \dS  = E_z(z,t)s = \frac{q_int(t)}{\epz} = \frac{\sigma(t)s}{\epz} \Rightarrow E_z(z,t) = \frac{\sigma(t)}{\epz}.
$$
Or $\sigma(t) = Q(t)/S$ et donc 
\eqbox{
  \Vec{E}_0(M,t) = \frac{Q(t)}{S\epz} \ez.
}
\end{liste}

}
\item Que vaut alors le vecteur densité de courant de déplacement $\jd$  ? Commenter la pertinence de l'ARQS magnétique. Déterminer le champ magnétique $\Vec{B}_i(M,t)$ induit par ce courant de déplacement de deux façons différentes : par une équation locale, et par un théorème intégral.
\sol{
  Par définition,
  $$
  \j_d = \epz \dr{\E}{t} = \frac{\dot{Q}(t)}{S} \ez.
  $$
Le courant de déplacement est non nul alors que le courant de conduction $\j$ est nul : on ne se situe pas dans le cadre de l'AQRS magnétique, mais celle de l'ARQS électrique ! On peut obtenir rapidement le champ magnétique engendré par ce courant de déplacement par l'équation de Maxwell-Ampère,
$$
\rot \B_i = \mu_0 \epz \dr{\E_0}{t},
$$
ou bien par le théorème d'Ampère généralisé. Les invariances et symétries de la distribution de courant de déplacement indique que $\Vec{B}_i(M,t) = B_i(z,t)\et$. On prend comme contour d'Ampère un cercle de rayon $r$ à l'altitude $z$ entre les armatures,
$$
\oint_C \B_i(M,t) \ps \d\vv{l} = B_i(z,t)2\pi r =  \mu_0 I_d = \mu_0 \frac{\dot{Q}(t)}{S} \pi r^2,
$$
et donc 
\eqbox{
  \Vec{B}_i(M,t) = \frac{\mu_0 \dot{Q}(t)}{2S} r \et.
}
}
\item Ce champ magnétique dépend-il du temps ? Quel phénomène physique intervient alors ? Déterminer le champ électrique induit $\Vec{E}_i(M,t)$ de deux façons différentes.
\sol{
  Le champ magnétique $\Vec{B}_i(M,t)$ dépend du temps. Le phénomène d'induction électromagnétique engendre un champ électrique induit $\Vec{E}_i(M,t)$, qui est donné par l'équation de Maxwell-Faraday,
  $$
  \rot \E_i = -\dr{\B_i}{t} = -\frac{\mu_0 \ddot{Q}(t)}{2S} r \et,
  $$
  ce qui s'intègre en 
  \eqbox{
    \Vec{E}_i(M,t) = -\frac{\mu_0 \ddot{Q}(t)}{4S} r^2 \ez.
  }

}
\item À quelle condition ce champ électrique induit est-il négligeable devant le champ électrique $\Vec{E}_0(M,t)$ pour $a=\SI{15}{cm}$ ?
\sol{
  On calcule le rapport des normes 
  $$
  \frac{E_i}{E_0} = \frac{\mu_0 \ddot{Q}(t)}{4S} r^2 \frac{S\epz}{Q(t)} = \frac{\mu_0 \ddot{Q}(t)}{4\epz Q(t)} r^2 \simeq \frac{\mu_0 \epz \omega^2}{4} a^2,
  $$
  où $\omega$ est la pulsation caractéristique des grandeurs électriques. Le champ électrique induit est négligeable devant le champ électrique $\Vec{E}_0(M,t)$ à basse fréquence, pour 
  \eqbox{
    \omega \ll \frac{2a}{\sqrt{\epz \mu_0}} \simeq \SI{e9}{rad.s^{-1}}.
  } 
}
\end{quest}
\textsc{Données :} Opérateur rotationnel en coordonnées cylindriques

 \begin{equation*}
  \rot{A} =  \left[\frac{1}{r} \frac{\partial A_{z}}{\partial \theta} - \frac{\partial A_{\theta}}{\partial z}\right]   \vec{e_r} +
				   \left[\frac{\partial A_{r}}{\partial z} - \frac{\partial A_{z}}{\partial r}\right]   \vec{e_\theta} +
				   \left[\frac{1}{r} \frac{\partial (rA_{\theta})}{\partial r} - \frac{1}{r} \frac{\partial A_{r}}{\partial \theta}\right]   \vec{e_z}
 \end{equation*}
\end{exo}
%%%%%%%%%%%%%%%%%%%%%%%%%%%%%%%%%%%%%%%%%%%%%%%%%%%%%%%

%%%%%%%%%%%%%%%%%%%%%%%%%%%%%%%%%%%%%%%%%%%%%%%%%%%%%%%
\begin{exo}[PC][2][colle]{Condensateur plan idéal en régime harmonique}
  Les armatures d’un condensateur plan sont deux disques de rayon $a$, distants de $e$, d’axe $(Oz)$.
  On applique une différence de potentiels sinusoïdale entre les armatures du condensateur.
   % $u(t)=u_0\cos(\omega t)$.
%\begin{figure}[h!]
%\centering
% \begin{center}
% \includegraphics[width=0.55\textwidth]{condo.pdf}
% \end{center}
%\end{figure}
Dans un premier temps, on suppose que le champ électrique est uniforme dans l’espace entre les armatures et s’écrit $\E_0 = E_0 \cos(\omega t)\ez$. 

\begin{quest}
\item Déterminer le champ magnétique $\B_1$ entre les armatures du condensateur.
\sol{
  Le champ magnétique $\B_1$ est induit par le champ électrique $\E_0$, et s'obtient soit par l'équation de Maxwell-Faraday (forme locale), soit par la loi de Faraday (forme intégrale). Sachant que $\E_0$ est porté par $\ez$, on s'attend à ce que $\B_1$ soit porté par $\et$ et dépende de $r$ uniquement, par analogie avec un fil infini parcouru par un courant. On suppose donc 
  % $$
  \eqbox{
    \B_1 = B_1(r,t) \et.
    }
L'équation de \udl{Maxwell-Ampère} donne alors 
$$\rot \B_1 =  \frac{1}{r}\dr{rB_\theta}{r} \ez = \mu_0\epz \dr{\E}{t} = -\omega\mu_0\epz E_0 \sin(\omega t)\ez,$$ ce qui donne
$$
B_1(r,t) = -\frac{r}{2} \omega\mu_0\epz E_0 \sin(\omega t) + \frac{f(t)}{r},
$$
où $f(t)$ est une fonction de temps. Or le champ ne peut pas diverger en $r=0$, ce qui impose $f(t)=0$. Donc
\eqbox{
\B_1 = -\frac{r}{2} \omega\mu_0\epz E_0 \sin(\omega t) \et.
} 
}
\end{quest}
Ce champ magnétique $\B_1$ étant variable, il est lui-même source d’un champ électrique $\E_2$. 

% On rappelle que $\mu_0\epz = 1/c^2$.

\begin{quest}
\item Déterminer $\E_2$ et montrer que le champ électrique approché s’écrit
$$
\E = \left[ 1-\frac{1}{4}\left(\frac{r\omega}{c}\right)^2\right]\E_0.
$$
À quelle condition se trouve-t-on dans le cadre de l’ARQS ?
Le champ électromagnétique $(\E,\B_1)$ est-il compatible avec l’équation de Maxwell-Ampère ?

\sol{
L’équation de \udl{Maxwell-Faraday} donne 
$$\rot \E_2 = -\dr{\B_1}{t} = \frac{r\omega^2}{2c^2}  E_0 \cos({\omega t}) \et.$$ Pour donner un rotationnel fonction de $r$ et porté $\et$, on doit choisir un champ fonction de $r$ et porté par $\ez$. On cherche donc $\E_2$ de la forme 
\eqbox{
\E_2 = E_2(r,t) \ez.
}
On trouve alors
$$
-\dr{E_2}{r} =  \frac{r\omega^2}{2c^2}  E_0 \cos({\omega t}) \Rightarrow E_2 = \frac{1}{4}\left(\frac{r \omega}{c}\right)^2 E_0 \cos({\omega t}) + g(t) .
$$
En l’absence de champ électrique uniforme $E_0$, on s’attend à ce qu’il n’y ait pas de champ électrique induit non plus, ce qui impose $g(t)=0$.
Donc
\eqbox{
  \E_2 = -\frac{1}{4}\left(\frac{r \omega}{c}\right)^2 \E_0.
  }
  En utilisant le principe de superposition, on trouve le champ électrique total, à cet ordre d’approximation,
\eqbox{
\E = \left[ 1-\frac{1}{4}\left(\frac{r\omega}{c}\right)^2\right]\E_0 .
}
}
\end{quest}

On cherche une solution exacte du problème sous la forme
$
\E(M,t) = E(r)\cos(\omega t)\ez.
$
% $$
% \E(r,t) = E(r)e^{i\omega t}\ez, \quad \text{avec} \quad E(r) = E_0\sum_{n=0}^{\infty} a_n \left(\frac{r \omega}{c} \right)^n.
% $$

\begin{quest}
  \item  Montrer que la fonction $E(r)$ vérifie l’équation 
  \begin{equation}
  \label{eq:condo}
  \frac{1}{r}\dd{}{r}\left(r \dd{E}{r} \right) = - k^2 E,
  \end{equation}
   avec $k=\frac{\omega}{c}$.
  %, et en déduire une relation de récurrence sur les coefficients $a_n$.
  % \textsc{Données} : $\Delta E(r) = \frac{1}{r}\dd{}{r}\left(r \dd{E}{r} \right)$.
  % \question
   On pose $E(r) = E_0 J(x)$ avec $x=r\omega/c$. On admet que le premier zéro de la fonction $J$ est en $x_0 \simeq \num{2.4}$. À partir de quelle fréquence observe-t-on une région de champ nul dans le condensateur pour $a = \SI{4}{cm}$ ?

   \sol{
On peut démontrer à partir des équations de Maxwell que le champ électrique est solution de l’équation de d’Alembert, soit 
$$
\Delta \E - \frac{1}{c^2} \ddr{\E}{t} = 0.
$$
En injectant la forme proposée pour $\E$, on trouve
$$
\Delta E(r) = \frac{1}{r}\dd{}{r}\left(r \dd{E}{r} \right) = -\frac{\omega^2}{c^2} E(r) = -k^2 E(r).
$$
On reconnaît l’équation de Helmholtz, qui admet des solutions de la forme $E(r) = E_0 J(x)$ avec $x=r\omega/c$. La fonction $f$ vérifie alors l’équation
$$
\frac{1}{x}\dd{}{x}\left(x \dd{J}{x} \right) = - J.
$$
On peut démontrer qu’une des solutions de cette équation est la fonctions de Bessel d'ordre 0, soit $J(x) = J_0(x)$, dont le premier zéro est en $x_0 \simeq \num{2.4}$. Le champ s’annule au bord du condensateur lorsque 
$$
\frac{a \omega}{c} = x_0 \Rightarrow \omega = \frac{x_0 c}{a} = \num{1.5e10} \text{ rad.s}^{-1},
$$
et donc 
\eqbox{
f = \SI{2.4e9}{ Hz}.
}

   }

\end{quest}
\end{exo}
%%%%%%%%%%%%%%%%%%%%%%%%%%%%%%%%%%%%%%%%%%%%%%%%%%%%%%%


%%%%%%%%%%%%%%%%%%%%%%%%%%%%%%%%%%%%%%%%%%%%%%%%%%%%%%%
\begin{exo}[PC][2][colle]{Condensateur plan idéal dans l'ARQS électrique}%{1}
L'étude du condensateur se traite rigoureusement dans le régime de l'ARQS électrique qui ne figure pas au programme ; nous admettrons que dans ce cas l'équation de Maxwell-Faraday est équivalente au cas statique, soit
\begin{equation}
  \label{eq:faraday}
  \rot{\vv{E}} = 0,
\end{equation}
toutes les autres équations de Maxwell étant inchangées.
On suppose dans la suite qu'un condensateur plan idéal est placé en série avec une résistance $R$ et un générateur de tension sinusoïdale $u(t)=u_0\cos(\omega t)$.
%\begin{figure}[h!]
%\centering
\begin{center}
\includegraphics[width=0.55\textwidth]{condo.pdf}
\end{center}
%\end{figure}
\begin{quest}
\item Donner l'expression du champ électrique $\vv{E}(M,t)$ entre les armatures du condensateur.
\item À partir de l'équation de Maxwell-Ampère,  déterminer l'expression du champ magnétique $\vv{B}(M,t)$.
\item Comparer les ordres de grandeur des densités d'énergie électrique $u_e$ et magnétique $u_m$. Conclure sur la nature de l'énergie prédominante stockée dans un condensateur.
\item Calculer le vecteur de Poynting $\Poy(M,t)$, puis la puissance $\mc{P}$ traversant les parois latérales du condensateur, de rayon $a$ et de hauteur $d$. Interpréter.
% \item Vérifier explicitement sur cet exemple le théorème de Poynting.
\end{quest}
\end{exo}
%%%%%%%%%%%%%%%%%%%%%%%%%%%%%%%%%%%%%%%%%%%%%%%%%%%%%%%


%%%%%%%%%%%%%%%%%%%%%%%%%%%%%%%%%%%%%%%%%%%%%%%%%%%%%%%%%%%%%%%%%%%%%%%%%%
\begin{exo}[PC][1][TD]{Énergétique d'un conducteur sphérique}
  Une boule chargée de centre $\rm O$ et de rayon $R$ se décharge dans un milieu infini ohmique de conductivité~$\sigma$. Le champ électromagnétique est nul à l'intérieur du conducteur. À l'extérieur, il s'écrit
  $$\vv{E}(M,t)=\frac{Q_0\exp(-\frac{\sigma t}{\varepsilon_0})}{4\pi\varepsilon_0 r^2}\,\vv{e_r}\qquad\text{et}\qquad \vv{B}(M,t)=0.$$
  \begin{quest}
    \item Expliquer sans exposer tous les calculs d'où provient l'expression donnée des champs.
    \item Calculer l'énergie électromagnétique à l'instant $t$.
    \item Calculer l'énergie dissipée par effet Joule dans le conducteur ohmique entre l'instant initial $t=0$ et un instant $t$ quelconque. Commenter.
  \end{quest}
\end{exo}
%%%%%%%%%%%%%%%%%%%%%%%%%%%%%%%%%%%%%%%%%%%%%%%%%%%%%%%%%

%%%%%%%%%%%%%%%%%%%%%%%%%%%%%%%%%%%%%%%%%%%%%%%%%%%%%%%%%%
\begin{exo}[PC][1][TD]{Champ électromagnétique sinusoïdal}
  \Notions{\item Équations de Maxwell \task Vecteur de Poynting}
  On considère le champ électrique suivant, régnant dans une partie de l’espace vide,
$$
\Vec{E}(M,t) = E_0 \cos(\omega t - k z) \ex + E_0 \sin(\omega t - k z) \ey.
$$

\begin{quest}
	\item Vérifier la compatibilité de cette expression avec l'équation Maxwell-Gauss.
	\sol{
  $$
  \Dive \E = \dr{E_x}{x}\ex + \dr{E_y}{y}\ey + \dr{E_z}{z} \ez = 0,
  $$
  compatible avec l'équation de Maxwell-Gauss \fbox{$\Dive \E = \frac{\rho}{\epz}$} car $\rho=0$ dans le vide.

	}
  \item Déterminer le champ magnétique $\B(M,t)$ associé. Quelle relation existe-t-il entre $k$ et $\omega$ ?
  \sol{
    L'équation de Maxwell-Faraday donne
    $$-\dr{\B}{t} = \rot\E = \left( \dr{E_y}{x}-\dr{E_x}{y} \right)\ez + \left( \dr{E_z}{y}-\dr{E_y}{z} \right)\ex + \left( \dr{E_x}{z}-\dr{E_z}{x} \right)\ey  = \dr{E_x}{z} \ey -\dr{E_y}{z}\ex $$
    Donc
    $ -\dr{\B}{t} = kE_0(\sin(\omega t + k z)\ey+\cos(\omega t + k z)\ex) = k \E $
    et donc en intégrant par rapport au temps,
    $$
    \B(M,t) = \frac{kE_0}{\omega} \left[ -\sin(\omega t + k z)\ex+\cos(\omega t + k z)\ey \right].    $$
    Avec $B_0 = \frac{kE_0}{\omega}$, on trouve 
    \eqbox{
    \B(M,t) = -B_0 \sin(\omega t - k z)\ex+B_0 \cos(\omega t - k z)\ey.
    }
    Les champs doivent aussi vérifier l’équation de Maxwell-Faraday, avec $\j = \Vec{0}$ dans le vide,
    \begin{align*}
      \rot\B & = \epz \mu_0 \dr{\E}{t} \\ 
      & \Rar kB_0\, \left[ \cos(\omega t - k z)\ey - \sin(\omega t - k z) \ex  \right] = \epz \mu_0 \omega E_0 \left[ \cos(\omega t - k z)\ey - \sin(\omega t - k z)\ex  \right] \\ 
      & \Rar k B_0 = -\epz \mu_0 E_0 \omega
    \end{align*}
    et donc 
    \eqbox{
    k = \pm \sqrt{\epz \mu_0} \omega.
    }
  }
	\item Déterminer le vecteur de Poynting $\Poy(M,t)$. 
  % Déterminer sa moyenne temporelle $\moy{\Poy}$ sur une période . 
  Représenter les vecteurs $\E$, $\B$ et $\Poy$ dans l’espace.
  \sol{
    \begin{align*}
      \Poy = \frac{1}{\mu_0}(\E \pv \B) & = \frac{E_0B_0}{\mu_0} (\cos(\omega t - k z) \ex + \sin(\omega t - k z) \ey) \pv ( -\sin(\omega t + k z)\ex +\cos(\omega t + k z)\ey) \\ 
      & = \frac{E_0B_0}{\mu_0} (\cos(\omega t + k z)^2 + \sin(\omega t + k z)^2) \ez.
    \end{align*}
      et donc 
      \eqbox{
        \Poy = \pm \sqrt{\frac{\epz}{\mu_0}} E_0^2 \ez
      }
      Les vecteurs $(\E,\B,\Poy)$ forment une triade directe. On rappelle que $\Pi$ représente la puissance surfacique transportée par le champ électromagnétique. La direction de propagation de l'énergie est donnée par $\pm \ez$ suivant le signe de $k$.
    }
  \end{quest}
  \end{exo}
  %%%%%%%%%%%%%%%%%%%%%%%%%%%%%%%%%%%%%%%%%%%%%%%%%%%%%%%%%
  
  
  %%%%%%%%%%%%%%%%%%%%%%%%%%%%%%%%%%%%%%%%%%%%%%%%%%%%%%%%%%
  \begin{exo}[PC][1][TD]{Champ électromagnétique atténué}
    \Notions{\item Équations de Maxwell \task Vecteur de Poynting}
    On considère le champ électrique suivant dans un milieu conducteur occupant l'espace $z>0$,
  $$
  \Vec{E}(M,t) = E_0 \cos(\omega t - k z) e^{-kz} \ex.
  $$
  
  \begin{quest}
    \item Que vaut la densité volumique de charge $\rho$ ?
    \sol{
    L'équation de Maxwell-Gauss donne
  $$
  \Dive \E = \dr{E_x}{x}\ex + \dr{E_y}{y}\ey + \dr{E_z}{z} \ez = 0 = \frac{\rho}{\epz},
  $$
  donc \fbox{$\rho=0$} dans un milieu conducteur.

	}
  \item Déterminer le champ magnétique $\B(M,t)$ associé. Que vaut le vecteur densité de courant $\j$ ?
  \sol{
  L'équation de Maxwell-Faraday donne
    $$-\dr{\B}{t} = \rot\E = \left( \dr{E_y}{x}-\dr{E_x}{y} \right)\ez + \left( \dr{E_z}{y}-\dr{E_y}{z} \right)\ex + \left( \dr{E_x}{z}-\dr{E_z}{x} \right)\ey  = \dr{E_x}{z} \ey $$
    Donc
    $ -\dr{\B}{t} = kE_0e^{-kz}\left[\sin(\omega t - k z)-\cos(\omega t - k z)\right] \ey $
    et donc en intégrant par rapport au temps, on trouve 
    \eqbox{
    \B(M,t) = B_0 e^{-kz} \left[\cos(\omega t - k z)+\sin(\omega t - k z) \right] \ey.
    }
    Le vecteur densité de courant s'obtient à partir de l’équation de Maxwell-Faraday,
    \begin{align*}
      \mu_0 \j  & = \rot\B - \epz \mu_0 \dr{\E}{t} \\ 
      & = \dots 
    \end{align*}
    Le courant est non nul dans un conducteur.
    }
	\item Déterminer le vecteur de Poynting $\Poy(M,t)$. Effectuer un bilan d'énergie sur un volume de conducteur.
  \sol{
    \begin{align*}
      \Poy = \frac{1}{\mu_0}(\E \pv \B) & 
    = \frac{E_0B_0}{\mu_0} \cos(\omega t - k z) e^{-kz} \ex \pv e^{-kz} \left[\cos(\omega t - k z)+\sin(\omega t - k z) \right] \ey \\ 
    & = \frac{E_0B_0}{\mu_0} e^{-2kz} \left[  \cos(\omega t - k z)^2 +  \cos(\omega t - k z)\sin(\omega t - k z) \right] \ez\, 
    \end{align*}
  On peut effectuer la moyenne temporelle sur une période pour obtenir la puissance moyenne. On utilise $\moy{\cos^2}=\frac{1}{2}$ et $\moy{\cos \sin} = 0$, soit
  \eqbox{
    \moy{\Poy} = \frac{k}{\omega \mu_0} \frac{E_0^2}{2} \ez
  }
  }
\end{quest}
\end{exo}
%%%%%%%%%%%%%%%%%%%%%%%%%%%%%%%%%%%%%%%%%%%%%%%%%%%%%%%%%


%%%%%%%%%%%%%%%%%%%%%%%%%%%%%%%%%%%%%%%%%%%%%%%%%%%%%%%%%%
\begin{exo}[PC][3][TD]{Câble coaxial en régime variable}
  \Notions{\item Équations de Maxwell \task Géométrie cylindrique \task Câble coaxial}

  {\begin{wrapfigure}{r}{0.25\textwidth}
    \vspace{-0.5cm} \centering \includegraphics[width=0.25\textwidth]{coax.png} \end{wrapfigure}
  Un câble coaxial est constitué de deux cylindres métalliques d’axe $(Oz)$, de rayons intérieur $R_1$ et extérieur $R_2 > R_1$. La couronne cylindrique située entre ces deux conducteurs a les propriétés électromagnétiques du vide.
Elle est le siège de champs électrique et magnétiques sinusoïdaux, dont on cherche l'expression en coordonnées cylindriques $(r,\theta,z)$ sous la forme :
\begin{equation}
  \label{eq:champ}
\Vec{E}(M,t) = E_0(r) \cos(\omega t - k z) \er, \qquad \B(M,t) = \Vec{B_0}(r) \cos(\omega t -k z).
\end{equation}
On admet que les champs $\Vec{E}(M,t)$ et $\Vec{B}(M,t)$ sont nuls à l'extérieur du câble, pour $r < R_1$
et $r > R_2$.}
\begin{quest}
	\item Déterminer la fonction $E_0(r)$ à partir d'une équation de Maxwell. On pourra poser $E_1 = E_0(R_1)$.
  % On posera $E_a = \lim\limits_{r \rightarrow a^+} E(r)$. En déduire le potentiel scalaire $V (r,z,t)$ qui sera pris nul en tout point du conducteur extérieur.
	\sol{
    L'équation de \udl{Maxwell-Gauss} dans le vide entre les conducteurs donne
    $$
    \Dive \E = \frac{1}{r}\dr{(rE_r)}{r} + \frac{1}{r} \dr{E_\theta}{\theta} + \dr{E_z}{z} = \frac{1}{r}\dd{(rE_0)}{r}  =  0.
    $$
    On intègre cette équation par rapport à $r$ pour obtenir
    $$
    \dd{(rE_0)}{r} = 0 \Rightarrow E_0(r) = \frac{A}{r},
    $$
    avec $E_0(R_1) = E_1 = \frac{A}{R_1}$. On en déduit 
    \eqbox{
      E_0(r) = E_1 \frac{R_1}{r}.
    }
	}
  \item Calculer le champ vectoriel $\Vec{B_0}(r)$. Simplifier son expression à partir de la relation de dispersion $\omega = k c$.
  \sol{
    On utilise l'équation de \udl{Maxwell-Faraday},
    $$
    \rot \E = -\dr{\B}{t} =   \dr{E_r}{z} \et  = E_0(r) k \sin(\omega t - k z) \et,
    $$
    soit 
    % \eqbox{
    $$
    \B(M,t) = E_0(r)  \frac{k}{\omega} \sin(\omega t - k z) \et.
    $$
    Or $\omega = k c$ d'après la relation de dispersion dans le vide, et donc 
    \eqbox{
    \Vec{B_0}(r) = \frac{E_a a}{rc} \et.
    }
    % }
  
  }
  % \item Quelle est la relation entre $k$ et $\omega$ ? On donne $\rot(\frac{\et}{r}) = \Vec{0}$.
  % \item Déteminer les densités surfaciques de charges $\sigma_a$ et $\sigma_b$ et de courant $\j_a$ et $\j_b$ sur les armatures cylindriques (en $r = a$ et $r = b$).
% Vérifier que l’équation de conservation de la charge sur les armatures est bien compatible avec
% les expressions trouvées.
\end{quest}
On note $U (z,t) = V_2(z,t) - V_1(z,t)$ la tension entre le conducteurs extérieur et intérieur, et $I (z,t)$ l’intensité du courant électrique porté par le conducteur intérieur.
\begin{quest}
  \item Déterminer l’expression de $U(z,t)$ et de $I(z,t)$.
  \sol{
    \begin{liste}
      \item \udl{Tension} : la tension est définit par la circulation du champ électrique entre les conducteurs, soit 
      $$
  U(z,t) = \int_a^b \Vec{E}(M,t) \ps \d r \,\er = \int_a^b E_a \frac{a}{r} \cos(\omega t - k z) \d r,
      $$
      soit 
      \eqbox{
        U(z,t) = E_a a \cos(\omega t - k z) \ln\left(\frac{b}{a}\right).
      }
      \item \udl{Intensité} : l'intensité peut s'obtenir à partir de \udl{théorème d'Ampère}. On choisit pour contour d'Ampère $\Gamma$ un cercle de rayon $r=a^+$ centré sur l'axe $(Oz)$. On trouve 
      $$
      \int_\Gamma \B(M,t) \ps \dl = 2\pi a \times \frac{E_a}{c} \cos(\omega t - k z) = \mu_0 I_\Gamma(z,t)  + \frac{1}{c^2} \dd{\Phi_E}{t} = \mu_0 I(z,t),
      $$
      car le courant de conduction enlacé par le contour vaut $I(z,t)$, et le flux du champ électrique est nul, puisque le champ est nul dans le conducteur, pour $r < a$. On en déduit 
      \eqbox{
      I(z,t) = \frac{2\pi a E_a}{\mu_0 c} \cos(\omega t - k z).
      }
    \end{liste}
  }
  \item Calculer l'impédance $Z$ du câble coaxial, puis vérifier que $Z$ est indépendant de $z$ et de $t$.
  \sol{
    L'\udl{impédance} $Z$ du câble coaxial est le rapport entre la tension et l'intensité,
    $$
    Z = \frac{U(z,t)}{I(z,t)} = \frac{E_a a \cos(\omega t - k z) \ln\left(\frac{b}{a}\right)}{\frac{2\pi a E_a}{\mu_0 c} \cos(\omega t - k z)},
    $$
    soit 
    \eqbox{
     Z = \frac{\mu_0 c}{2\pi} \ln\left(\frac{b}{a}\right).
    }
  L'impédance est indépendante de $z$ et de $t$, et ne dépend que des propriétés géométriques du câble coaxial.
  }
  \item Déterminer le vecteur de Poynting puis la puissance moyenne qui traverse une section droite
du câble. Montrer qu’elle s’exprime simplement en fonction de $Z$ et commenter.
\sol{
  \begin{liste}
    \item \udl{Vecteur de Poynting} : on calcule le vecteur de Poynting par la relation 
    $$
    \Poy(M,t) = \frac{1}{\mu_0}(\E \pv \B) = \frac{E_0(r) {B_0}(r)}{\mu_0} \cos^2(\omega t - k z) \ez,
    $$
    soit 
    \eqbox{
      \Poy(M,t) = \frac{E_a^2 a^2}{\mu_0 c r^2} \cos^2(\omega t - k z) \ez.
    }
    La moyenne temporelle du vecteur de Poynting donne la puissance moyenne surfacique transportée par le champ électromagnétique,
    $$
    \moy{\Poy} = \frac{E_a^2 a^2}{2\mu_0 c r^2} \ez.
    $$
    \item \udl{Puissance moyenne} : la puissance moyenne qui traverse une section droite du câble coaxial de rayon $r$ est donnée par l'intégrale de la puissance surfacique sur la surface de la section droite,
    $$
    \mc{P} = \iint_S \moy{\Poy}\ps \dS = \int_a^b \int_0^{2\pi} \frac{E_a^2 a^2}{2\mu_0 c r^2} r\d r \d \theta = \frac{E_a^2 a^2}{2\mu_0 c} \int_0^{2\pi} \d\theta \int_0^{2\pi r} \frac{\d r}{r}  = \frac{\pi E_a^2 a^2}{\mu  _0 c} \ln\left( \frac{b}{a} \right).
    $$
    On peut aussi écrire la puissance moyenne en fonction de l'impédance du câble coaxial,
    \eqbox{
      \mc{P} = \frac{1}{2} Z I_0^2,
    }
    avec 
    $$
    I_0 = \frac{2\pi a E_a}{\mu_0 c}$$
    l'amplitude du courant électrique. On retrouve la loi de Joule en régime sinusoïdal.
  \end{liste}
}
\end{quest}

\textsc{Données :} 
% \begin{tabular}{l}
 $\rot (\lambda \Vec{a}) = \lambda \,\rot \Vec{a} + \grad \lambda \pv \Vec{a}$,
  $\rot(f(r)\,\er) = \Vec{0}$,
  $\rot\left( \textstyle\frac{\et}{r} \right) = \Vec{0}$.
% \end{tabular}

\end{exo}
%%%%%%%%%%%%%%%%%%%%%%%%%%%%%%%%%%%%%%%%%%%%%%%%%%%%%%%%%

